\chapter{Glossary and reference}
\label{ch:glossary}

\begin{aside}
Part II has introduced a lot of vocabulary. This chapter is a
reference: every term in one place, with a one-paragraph
explanation and a cross-reference back to its detailed chapter.
Use it as a refresher before tackling Parts III and IV.
\end{aside}

\section{Architecture terms}

\textbf{Element.} The unit of UI tree. A sum of
\code{BoxElement}, \code{TextElement}, \code{ElementList}.
Chapter~\ref{ch:elements}.

\textbf{Model.} The application state. A struct holding every
piece of state the app needs. Pure data; no pointers, no
shared state. Chapter~\ref{ch:elm}.

\textbf{Msg.} The alphabet of state transitions. A
\code{std::variant} where each alternative is one kind of
event the program responds to. Chapter~\ref{ch:elm}.

\textbf{Cmd<Msg>.} A description of side effects. A sum type
whose alternatives include \code{None}, \code{Quit},
\code{Task}, \code{After}, etc. Interpreted by the runtime.
Chapter~\ref{ch:elm}.

\textbf{Sub<Msg>.} A description of event sources the program
listens to. A sum type whose alternatives include \code{OnKey},
\code{OnMouse}, \code{Every}, etc. Chapter~\ref{ch:elm}.

\textbf{view.} A pure function from \code{Model} to
\code{Element}. Called every frame. Chapter~\ref{ch:elm}.

\textbf{update.} A pure function from
$(\text{Model}, \text{Msg})$ to $(\text{Model}, \text{Cmd})$.
Drives state transitions. Chapter~\ref{ch:elm}.

\textbf{subscribe.} A pure function from \code{Model} to
\code{Sub}. Declares which event sources are active.
Chapter~\ref{ch:elm}.

\textbf{init.} A pure function returning the initial
\code{Model} and \code{Cmd}. Called once on startup.
Chapter~\ref{ch:elm}.

\textbf{Program.} A concept aggregating the four functions
and their types. \code{maya::run} takes a \code{Program}.
Chapter~\ref{ch:concepts}.

\section{Rendering terms}

\textbf{Pipeline.} The six-stage transformation from
\code{Element} to ANSI bytes: view, layout, paint, diff,
serialize, write. Chapter~\ref{ch:pipeline}.

\textbf{Layout.} The phase that assigns each node a final
\code{Rect}. Uses a flexbox-style algorithm with integer cells.
Chapter~\ref{ch:layout}.

\textbf{Canvas.} A flat grid of packed cells, sized to the
terminal. Painted from the layout result. Chapter~\ref{ch:pipeline}.

\textbf{Packed cell.} A 64-bit struct holding a codepoint, a
style ID, and flags. Chapter~\ref{ch:pipeline}.

\textbf{Diff.} The phase that compares the current canvas to
the previous one and produces a per-row dirty range. SIMD-
accelerated. Chapter~\ref{ch:pipeline}.

\textbf{Serialize.} The phase that turns the diff into a
minimal sequence of ANSI bytes: cursor moves, SGR sequences,
cell writes. Chapter~\ref{ch:pipeline}.

\textbf{Style pool.} An intern table mapping unique
\code{Style} structs to small integer IDs. Used to keep packed
cells small. Chapter~\ref{ch:pipeline}.

\textbf{Cache fingerprint.} A 64-bit hash of the element
tree. Equal fingerprints mean the renderer can skip the entire
pipeline. Chapter~\ref{ch:pipeline}.

\textbf{Force redraw.} A \code{Cmd} alternative that discards
the cache and emits every cell on the next frame. Used after
terminal disturbances. Chapter~\ref{ch:elm}.

\section{DSL terms}

\textbf{NTTP.} Non-Type Template Parameter. A value (not a
type) that appears in a template parameter list. C++20
relaxed the rules so structural types (\code{Str},
\code{CTStyle}) can be NTTPs. Chapter~\ref{ch:dsl}.

\textbf{consteval.} A function that must run at compile time.
The compiler errors on runtime invocation. The DSL's pipe
operators are typically \code{consteval}. Chapter~\ref{ch:constexpr-deep}.

\textbf{constexpr.} A function that may run at compile time if
its arguments are constant expressions. May also run at runtime.
Chapter~\ref{ch:constexpr-deep}.

\textbf{Str.} A structural type wrapping a fixed-size character
array. Lets string literals appear as NTTPs.
Chapter~\ref{ch:dsl}.

\textbf{CTStyle.} A compile-time style struct holding boolean
flags and colour fields. Forms the payload of \code{StyTag}.
Chapter~\ref{ch:dsl}.

\textbf{StyTag.} A wrapper template indexed by a \code{CTStyle}
NTTP. Two \code{StyTag} values with the same \code{CTStyle}
have the same type. Chapter~\ref{ch:dsl}.

\textbf{TextNode.} A compile-time element wrapping a literal
string and an optional style. Pipe operators on it remain
compile-time. Chapter~\ref{ch:dsl}.

\textbf{RuntimeTextNode.} The runtime counterpart of
\code{TextNode}. Holds an owning string and style.
Chapter~\ref{ch:dsl}.

\textbf{BoxNode.} A compile-time element representing a box
with children. Indexed by a \code{BoxCfg} NTTP.
Chapter~\ref{ch:dsl}.

\textbf{Type-state machine.} The pattern where the type of a
DSL node encodes which operations are legal on it. The pipe
overloads have \code{requires} clauses that enforce the
transitions. Chapter~\ref{ch:dsl}.

\textbf{Concept.} A named compile-time predicate on a type. Used
to constrain template parameters. Chapter~\ref{ch:concepts}.

\section{Effect and event terms}

\textbf{Task.} A \code{Cmd} alternative that runs a callable on
a worker thread; its return value becomes a message.
Chapter~\ref{ch:elm}.

\textbf{IsolatedTask.} Like \code{Task} but on a dedicated
thread, not the pool. For long blocking work.
Chapter~\ref{ch:elm}.

\textbf{Batch.} A \code{Cmd} or \code{Sub} alternative that
groups multiple effects or sources. Chapter~\ref{ch:elm}.

\textbf{After.} A \code{Cmd} alternative that delivers a
message after a duration. Used for debounce, retry backoff,
toast expiry. Chapter~\ref{ch:elm}.

\textbf{Every.} A \code{Sub} alternative that delivers a
periodic message. Used for animations and polling.
Chapter~\ref{ch:elm}.

\textbf{AnimationFrame.} A \code{Sub} alternative that delivers
a message after each render frame. Used for frame-synchronised
animation. Chapter~\ref{ch:elm}.

\textbf{OnKey, OnMouse, OnResize, OnPaste.} \code{Sub}
alternatives for the various input events. Chapter~\ref{ch:elm}.

\textbf{Cmd::map.} A function that lifts a
\code{Cmd<ChildMsg>} into a \code{Cmd<ParentMsg>} via a
mapping function. The mechanism for composing components.
Chapter~\ref{ch:elm}.

\section{Style and theme terms}

\textbf{Style.} A struct holding fg, bg, attributes, and
optional hyperlink. The per-cell colouring directive.
Chapter~\ref{ch:style}.

\textbf{Attribute.} A bitset of typographic flags (Bold, Dim,
Italic, etc.). Combine with bitwise OR. Chapter~\ref{ch:style}.

\textbf{Color.} A sum type representing default, Ansi16,
Ansi256, or Rgb. Quantises if the terminal does not support
the requested form. Chapter~\ref{ch:elements}.

\textbf{Theme.} A palette plus metadata. Maps named roles to
colours. Swappable at runtime. Chapter~\ref{ch:style}.

\textbf{Palette.} An array of \code{Color} values indexed by
\code{Role}. The application-level colour vocabulary.
Chapter~\ref{ch:style}.

\textbf{Role.} A named slot in the palette: Foreground,
Background, Primary, Accent, Warning, etc. Decouples ``what
the colour means'' from ``what the colour is''.
Chapter~\ref{ch:style}.

\section{Layout terms}

\textbf{FlexStyle.} The layout parameters on a box: direction,
justify, align, grow, shrink, basis, wrap, gap, padding,
border. Chapter~\ref{ch:layout}.

\textbf{Direction.} Row or Column. Determines main axis.
Chapter~\ref{ch:layout}.

\textbf{Justify.} Where children sit along the main axis:
Start, End, Center, SpaceBetween, SpaceAround, SpaceEvenly.
Chapter~\ref{ch:layout}.

\textbf{Align.} Where children sit along the cross axis:
Start, End, Center, Stretch, Baseline. Chapter~\ref{ch:layout}.

\textbf{Grow.} A non-negative number; children with positive
grow share extra space proportionally. Chapter~\ref{ch:layout}.

\textbf{Shrink.} A non-negative number; children with positive
shrink absorb deficit proportionally. Chapter~\ref{ch:layout}.

\textbf{Basis.} The starting size before grow/shrink.
Chapter~\ref{ch:layout}.

\textbf{Wrap.} What to do when children overflow: NoWrap, Wrap,
WrapReverse. Chapter~\ref{ch:layout}.

\textbf{Gap.} Fixed cell spacing between children.
Chapter~\ref{ch:layout}.

\textbf{Padding.} Per-side integer cells, inside the box.
Chapter~\ref{ch:layout}.

\textbf{Border.} A decoration around a box; takes one cell per
side. Chapter~\ref{ch:elements}.

\section{Runtime terms}

\textbf{Context.} A bundle of runtime resources (platform
handles, config, caps, logger, stats) passed to subsystems.
Chapter~\ref{ch:runtime}.

\textbf{EventLoop.} The driver that waits for events and drives
the pipeline. Chapter~\ref{ch:runtime}.

\textbf{EventSource.} The platform abstraction for the
underlying OS event mechanism (epoll, kqueue,
WaitForMultipleObjects). Chapter~\ref{ch:platform}.

\textbf{SubscriptionManager.} The component that diffs the
declared subscription set against the active OS-level wires.
Chapter~\ref{ch:runtime}.

\textbf{Interpreter.} The component that walks
\code{Cmd<Msg>} variants and performs the side effects.
Chapter~\ref{ch:runtime}.

\textbf{Renderer.} The component that runs the six-stage
pipeline and writes to stdout. Chapter~\ref{ch:runtime}.

\textbf{Wake fd.} A self-pipe used by worker threads to wake
the main loop. Chapter~\ref{ch:platform}.

\section{Reactivity terms}

\textbf{Signal<T>.} A value that notifies subscribers when it
changes. Chapter~\ref{ch:reactivity-preview}.

\textbf{Computed<T>.} A derived signal whose value is a
function of other signals. Recomputes on input change.
Chapter~\ref{ch:reactivity-preview}.

\textbf{Effect.} A callable that re-runs whenever any signal it
reads changes. Chapter~\ref{ch:reactivity-preview}.

\section{Other terms}

\textbf{Expected<T, E>.} A value-typed sum of success and
error. \maya{}'s value-typed error handling type.
Chapter~\ref{ch:expected}.

\textbf{IoError.} A struct describing an I/O failure: code,
context, syscall name. The default error type for boundary
operations. Chapter~\ref{ch:expected}.

\textbf{Error boundary.} A wrapper that catches exceptions in
\code{build}/\code{view} and renders a fallback element.
Chapter~\ref{ch:errorboundary}.

\textbf{FocusManager.} A small helper that tracks focus IDs
and order. Stored in the model.
Chapter~\ref{ch:focus}.

\textbf{ScrollState.} A small struct (offset, viewport,
total) representing scroll position. Used for virtualisation.
Chapter~\ref{ch:scroll}.

\textbf{GraphemeWalker.} An iterator over Unicode grapheme
clusters. Cell-width-aware. Chapter~\ref{ch:text}.

\textbf{Widget.} A reusable UI component. May be a pure
renderer, a stateful component, or a reactive widget.
Chapter~\ref{ch:widget-arch}.

\textbf{Z-stack.} A composition mode where children render at
the same position; the last one wins on conflict. Used for
modals, overlays. Chapter~\ref{ch:end-to-end}.

\textbf{Hot reload.} Restarting an app while preserving its
model. Chapter~\ref{ch:lifecycle}.

\textbf{NO\_COLOR.} An environment variable that, when set,
asks programs to disable colour. \maya{} honours it via the
monochrome theme. Chapter~\ref{ch:a11y}.

\textbf{OSC 8.} The ANSI escape for clickable hyperlinks.
\maya{}'s \code{TextElement::hyperlink} emits it.
Chapter~\ref{ch:elements}.

\textbf{Kitty Keyboard Protocol.} An extended keyboard reporting
mode supported by modern terminals. Provides better key
disambiguation. Foundations chapter (Part I).

\textbf{Bracketed paste.} A terminal mode where pasted text is
wrapped in escape sequences, distinguishable from typed input.
Chapter~\ref{ch:elm}.

\textbf{Inline mode.} A rendering mode where \maya{} draws to
the existing terminal scrollback without taking over the screen.
Treated fully in Chapter~\ref{ch:inline}.

\section{Index of cross-references}

Each chapter has a \code{\textbackslash label}; cross-references
use \code{\textbackslash ref}. To find where a term is treated
in depth, look up the term here and follow the chapter link.

For online reading, the PDF is hyperlinked; clicking a chapter
reference jumps to the target. For print reading, the page
numbers in the TOC are the navigation tool.

\section*{Workshop}

\begin{enumerate}[leftmargin=*]
  \item Pick five terms from this glossary you cannot define
    from memory. Re-read the relevant chapter sections.
  \item Map each term to its position in the architecture
    diagram on page xxx (the architecture overview).
  \item For each \code{Cmd} alternative listed in the glossary,
    write the smallest possible app that uses it.
\end{enumerate}

\begin{recap}
The glossary collects every term Part II introduced. Use it as
a reference when reading later parts. Three sentences from the
previous chapter summarise the architecture; this chapter is
the vocabulary supporting those sentences.
\end{recap}
